% Derived one-step double-sieve fragment, 2026-10-08.
% Not inserted in the main manuscript; no global Chen constant asserted.
% The full Wu H-table migration was rejected during independent review.

\subsection{A same-new-level one-step correction on a Buchstab block}
\label{sec:same-new-level-double-sieve-route}

Write $\mathcal A_N=\{N-P:P<N,\ P\text{ prime}\}$,
$K=681/100$, $\mu=25/128$, and $v(t)=(1+t)/2$ for $t<\mu$.
We use
\href{https://www.impan.pl/shop/en/publication/transaction/download/product/82961}
{Wu 2004, Lemma 4.1 and equations (5.1)--(5.10)},
\href{https://arxiv.org/html/2109.02851v2}
{Lichtman, arXiv:2109.02851v2, Lemma 6.4},
and the already derived moving-residue prefix interface based on
\href{https://arxiv.org/html/2309.08522v1}
{Lichtman, arXiv:2309.08522v1, Propositions 5.2 and 6.1}.
The $K$ input is the previously certified positive prime-product upper
bound from \texttt{extended-prefix-upper.tex}. The exact small
calculation below is saved in
\texttt{certify-same-new-level-double-sieve-route.py} and
\texttt{same-new-level-double-sieve-route.json}.

This construction does not infer the signed switched-distribution
hypothesis (5.22) of L21 from a bound for the number of primes. Instead
we use the positive switched count at its actual place in Wu's first
weighted inequality. We do not use the $H,h$ tables or their full
iteration. The resulting explicit function $J$ is a new-argument
correction to a linear upper bound at the same new total level.

\begin{proposition}[A one-step correction for the $L_2$ Buchstab terms]
\label{prop:same-new-level-one-step}
Let
\[
 \frac{31}{200}\le t\le\frac4{25},\qquad
 p=N^{t+o(1)},\qquad \nu=v(t)=\frac{1+t}{2},\qquad
 s=\frac{\nu-t}{t}=\frac{1-t}{2t}.
\]
On every positive averaged prefix box in this interval, after the
usual fixed level buffers and exponent-box limits, the permissible
upper multiplier for $S(\mathcal A_{N,p},p)$ is
\begin{equation}\label{eq:same-new-one-step-upper}
 F(s)-\frac{2e^\gamma}{s}J_0,\qquad J_0=\frac1{250}.
\end{equation}
It is a bound averaged over the prefix primes, with their usual
positive main-factor measure, not an estimate for each individual
prime. In the existing variable-level Buchstab lower bound for $L_2$,
this yields an additional lower coefficient
\begin{equation}\label{eq:same-new-one-step-L2}
 \Delta L_2\ge8J_0\log\frac{676}{651}.
\end{equation}
Consequently the complete Chen functional increases by at least
\begin{equation}\label{eq:same-new-one-step-Chen}
 \Delta c\ge2J_0\log\frac{676}{651}>
 \frac2{6635}.
\end{equation}
The last inequality is purely rational after using
$\log(676/651)=2\operatorname{atanh}(25/1327)>50/1327$.
\end{proposition}

\begin{proof}
Fix a compact exponent box $t\in[t_-,t_+]$ strictly inside
$(31/200,4/25)$.
Use a single total level $Q=N^\nu$ with
$\nu<v(t_-)$ and a fixed positive margin. Then choose all sieve,
factorization, local-scale and mesh precisions inside that margin.
The centre notation below denotes the subsequent shrinking-box limit.
The local exponents are
\[
 \nu\in[231/400,29/50],\qquad
 R=Q/p=N^{L+o(1)},\qquad L=\nu-t=\frac{1-t}{2},
\]
and
\[
 \frac{21}{8}\le s\le\frac{169}{62},\qquad s_0=\frac72,\qquad
 y=R^{1/s_0},\qquad z=p=R^{1/s}.
\]
All inequalities used below have strict surplus, so fixed buffers
preserve them.

For $\sigma(p)$ equal to the anchor prime-box indicator, or a fixed
finite nonnegative constant combination of such indicators, define
\[
 \Phi_s=\sum_p\sigma(p)S(\mathcal A_{N,p},p),\qquad
 \Theta_Q=4\operatorname{li}(N)
   \sum_p\frac{\sigma(p)\mathcal C(pN)}
                 {\varphi(p)\log(Q/p)}.
\]
The ordinary normalized main multipliers are
$A(s)=sF(s)/(2e^\gamma)$ and $a(s)=sf(s)/(2e^\gamma)$.
In particular $A(s)=1$ on this block.

Wu's first weighted inequality has the form
\begin{equation}\label{eq:same-new-one-step-weight}
 2\Phi_s\le
  2\sum_p\sigma(p)S(\mathcal A_{N,p},y)
 -\sum_p\sigma(p)\sum_{y\le q<z}S(\mathcal A_{N,pq},y)
 +\sum_p\sigma(p)\!\sum_{y\le q_1<q_2<q_3<z}
        S(\mathcal A_{N,pq_1q_2q_3},q_2)
 +o(\Theta_Q).
\end{equation}
The prime set excludes divisors of $pN$, with the same additional
prefix exclusions as in Wu's formula. Repeated-prime exceptions can
be removed with a power saving. The inequality follows from the
pointwise Buchstab identities in Wu (4.1)--(4.5); it does not require
the original short-box class after this fixed positive box has been
specified. Thus no error relative to a vanishing short-box mass is
being inferred from a fixed-mass asymptotic.

\paragraph{The two non-switched terms.}
The first term has the ordinary upper multiplier $2A(s_0)$.
Its residual primes are below $y<p$, so the largest modulus prime
remains the anchor. For the second term write
$q=N^{Lu+o(1)}$, $1/s_0\le u\le1/s$. Then
\[
 N^{3/25+o(1)}\le q\le p,\qquad
 t+2Lu\le3t\le\frac{12}{25}<\nu,\qquad
 2t\le\frac8{25}<\nu.
\]
The endpoint equality $q=p$ is trimmed before taking limits.
The lower residual argument is
\[
 s_0(1-u)\in[13/6,5/2].
\]
These are exactly the ordered two-prefix configurations covered by
the existing residual prime-prefix interface: the cutoff $y$ is at
most the smaller specified prime $q$, both initial square conditions
hold with fixed margins, and every later rough square condition
follows from its residual level. Empty residual rough atoms are
included. Both signs retain their original Iwaniec atom coefficients,
small-prime weights, and multiplicity factorials. Every concatenated
rough atom has largest exponent in $[t_-,t_+]$, and hence cap
$v(t_-)> \nu$; L23 Proposition 6.1 supplies its factorization.
The moving-residue argument uses Proposition 5.2 on the true local
scales, as already proved in \texttt{prefix-sieve-interface.tex}.

The negative term in \eqref{eq:same-new-one-step-weight} is therefore
bounded below, after division by $\Theta_Q$, by
\[
 B(s)=\int_{1/s_0}^{1/s}
  \frac{\log(s_0(1-u)-1)}{u(1-u)}\,du
 =\int_{1-1/s}^{1-1/s_0}
  \frac{\log(s_0x-1)}{x(1-x)}\,dx.
\]
No lower double-sieve correction is used here.

\paragraph{The exact positive switched mass.}
The last term counts primes $N-m$ in a labelled product multiset
\[
 m=pq_1q_2q_3n\le N,\qquad
 y\le q_1<q_2<q_3<z,\qquad P^-(n)\ge q_2,
\]
after the same coprimality and repeated-prime exceptions as in
Wu (5.3)--(5.4). Canonically factor $n$ into its ordered prime factors.
Every prime factor of $m$ is at least $N^{3/25+o(1)}$, which is
strictly larger than the $N^{100/1149}$ cutoff of the existing
$K$ theorem. The fixed level buffers are chosen small enough to
preserve this strict gap; the exact exponent $3/25$ is the unbuffered
limiting value. There are only finitely many prime-factor patterns.
All ordering, roughness, product, and prefix-box constraints become
finite monomial inequalities; their labels are retained. Thus the
existing positive-count theorem applies to this exact multiset.

The quotient $n=1$ contributes at most $N^{4t+o(1)}\le N^{16/25+o(1)}$
to the unsifted mass and to the desired count. It is harmless here
by this explicit power bound. Also, throughout the triple integration,
\[
 \alpha=\frac{\log(N/p)}{\log(Q/p)}\ge2+o(1),\qquad
 \frac{\alpha-u_1-u_2-u_3}{u_2}\ge2s-3\ge\frac94.
\]
The standard rough-number estimate is therefore used away from its
argument-one boundary. In particular no six-prime unit-quotient
boundary is present in this one-step construction.

Let $X$ be the unsifted labelled mass. The same partial summation as
in Wu (5.8) gives
\[
 X\le (I(s)+o(1))N
       \sum_p\frac{\sigma(p)}{\varphi(p)\log(Q/p)},
\]
where
\[
 I(s)=\sup_{\alpha\ge2}
 \int_{1/s_0\le u_1\le u_2\le u_3\le1/s}
 \omega\!\left(\frac{\alpha-u_1-u_2-u_3}{u_2}\right)
 \frac{du_1\,du_2\,du_3}{u_1u_2^2u_3}.
\]
For completeness $\omega(w)\le3/5$ for $w\ge2$ follows from
$\omega(w)=(1+\log(w-1))/w$ on $[2,3]$ and the delay equation.
The maximum of
$1+\log(w-1)-(3/5)w$ on $[2,3]$ is at $w=8/3$ and is negative:
$\log(5/3)\le((5/3)-(3/5))/2=8/15<3/5$.
The delay equation propagates this upper bound on successive unit
intervals. Hence
\begin{equation}\label{eq:same-new-one-step-Ibar}
 I(s)\le\overline I(s)=\frac35
 \left\{(s_0+s)\log\frac{s_0}{s}-2(s_0-s)\right\}.
\end{equation}
The elementary integration here first integrates $u_1,u_3$, leaving
$\int_{1/s_0}^{1/s}\log(s_0u)\log(1/(su))u^{-2}\,du$.

The $K$ upper bound, in its common-mass normalization, is
$(K+o(1))\mathcal C(N)X/\log N$. Since
$\mathcal C(N)\le\mathcal C(pN)$ and $\operatorname{li}(N)\sim N/\log N$,
the last term in \eqref{eq:same-new-one-step-weight}, divided by
$\Theta_Q$, is at most $(K/4)\overline I(s)+o(1)$.
The factor is $K/4$, not $K/(4/(\nu-t))$.
This is an upper bound for a positive prime count, not a signed
distribution assertion for a specified global sieve weight.
The $K$ proof's one common positive smooth outer approximation may be
used here on each fixed exponent box, followed by its smoothing limit;
there is no need to claim uniformity over Wu's shrinking short boxes.

\paragraph{An explicit new-argument saving.}
Since
\[
 A(7/2)=1+C,\qquad
 C=\int_2^{5/2}\frac{\log(x-1)}x\,dx,
\]
the three estimates in \eqref{eq:same-new-one-step-weight} give
\begin{equation}\label{eq:same-new-one-step-J}
 \Phi_s\le\{1-J(s)+o(1)\}\Theta_Q,\qquad
 J(s)=-C+\frac12B(s)-\frac{681}{800}\overline I(s).
\end{equation}
Every argument and every baseline in this formula uses the same $Q$.
No saving evaluated at $(1/2-t)/t$ is added.

Here is a short outward certificate for $J(s)>1/250$ on the entire
interval $[s_-,s_+]=[21/8,169/62]$. The function $B$ decreases because
its positive-integrand interval shrinks with $s$. Also
\[
 \overline I'(s)=\frac35\left(\log(s_0/s)+1-s_0/s\right)\le0.
\]
Consequently
\[
 J(s)\ge-C+\frac12B(s_+)-\frac{681}{800}\overline I(s_-).
\]
The integrands defining $C$ and $B(s_+)$ are increasing on their
integration intervals. For $C$ this follows from
$x/(x-1)>\log(x-1)$; for $B$, the positive numerator increases and
$x(1-x)$ decreases because $x>1/2$.
Use a right Riemann sum for $C$ and a left Riemann sum for $B(s_+)$,
each with $32768$ equal rational subintervals.
The existing $80$-bit integer interval logarithm routine, with its
explicit atanh remainder, gives the exact lower bound
\[
 -C_{\rm right}+\frac12B_{\rm left}
       -\frac{681}{800}\overline I(s_-)
 \ge
 \frac{4850531719234607697725}{1208925819614629174706176}
 >\frac1{250}.
\]
No floating point number decides this inequality. The saved certificate
also records a stronger optional check, using $16$ rational
$s$-subintervals and $8192$ cells in each one-dimensional Riemann sum;
its minimum lower bound exceeds $1/200$. The proposition only uses
the simpler bound $1/250$.

In terms of the ordinary main factor at cutoff $p$, the normalized
defect in \eqref{eq:same-new-one-step-J} is
$2e^\gamma J_0/s$, proving \eqref{eq:same-new-one-step-upper}.
The existing Buchstab lower bound for $L_2$ subtracts the ordinary
upper kernels for $S(\mathcal A_{N,p},p)$ over
$25/224<t<50/309$. Our entire block lies strictly inside this range.
Its ordinary kernel uses precisely the level $v(t)$, so the saving
can be inserted in that identity:
\[
 \Delta L_2\ge
 4J_0\int_{31/200}^{4/25}\frac{dt}{t(v(t)-t)}
 =8J_0\log\frac{(4/25)(169/200)}{(31/200)(21/25)}
 =8J_0\log\frac{676}{651}.
\]
The $L_2$ coefficient in the complete eleven-term weight is $1/4$.
This proves \eqref{eq:same-new-one-step-Chen}.
All parameters are fixed before $N\to\infty$; next remove the
smoothing and sieve-precision errors, shrink the boxes and their
level buffers within the open anchor interval, and only then exhaust
that interval by compact subintervals. On each such compact interval
the buffered arguments remain inside the certified $s$ interval.
The endpoints have zero main-factor mass. The strict numerical
surplus above $J_0$ absorbs the remaining errors.
\end{proof}

\paragraph{Why the full old $H$ table was not imported.}
The preceding proposition proves a positive correction directly.
It does not prove the stronger proposed statement that the entire
old table, such as $H(3)\ge72943/10^7$, transfers to all new-level
states. Independent review found a concrete interface gap in that
proposed extension.

In Wu's Lemma 4.2, terms $\Gamma_1$ and $\Gamma_5,\ldots,\Gamma_8$
are positive non-switched upper terms; $\Gamma_2,\Gamma_3,\Gamma_4$
are negative lower terms; and $\Gamma_9,\ldots,\Gamma_{21}$ are
positive switched terms. The latter consist of respectively three
additional primes for indices $9$--$15$, four for $16$--$19$, five
for $20$, and six for $21$. Replacing their final positive prime-count
upper bounds by $K$ would have the correct $K/4$ normalization.
That replacement alone does not check terms $5$--$8$.

Specifically, the published seed at $s=2.2$ uses $\phi_2=2.90$.
The $\Gamma_5$ domain allows two new exponents $u<v<1/\phi_2$,
measured relative to $R=Q/b$. Near its upper corner,
$v+2u>1$. The enlarged prefix then has square exponent
$T+L(v+2u)>\nu$, and is outside the full
$\mathcal D^{\rm well}(Q)$ class of the presently proved prefix
interface. Its small bare product does not solve this: residual
modulus atoms can raise the full modulus above $N^{1/2}$.
A new factorization theorem for that explicit enlarged-prefix domain,
or a replacement upper kernel and a recomputed recursion, is needed
before importing the old $\Psi_2$ seeds or their nine-by-nine matrix.

There is a second mass check to retain for such a future construction.
In the five- and six-extra-prime terms, the unit quotient may lie on
a product boundary at scale $N$. Keeping it in the exact multiset
does not by itself show that an integral involving only the ordinary
Buchstab function includes that boundary mass. Its negligibility,
or its separate mass term, must be checked before reusing those
$I_{2,i}$ envelopes. Neither issue occurs for the one-step triple
term proved above.

Thus the direct L21 black-box route still lacks its requested signed
switched-distribution statement at level $v(t)$. The proposition gives
a different, explicitly proved positive-count route on a concrete
$L_2$ block. It does not restore a universal mixed-level bridge and
does not independently certify a new global counting coefficient.
